\documentclass[12pt]{article}
\usepackage[utf8]{inputenc}
\usepackage[russian]{babel}
\usepackage{amsmath}
\usepackage{amssymb}
\usepackage{geometry}
\geometry{a4paper, margin=2cm}

\begin{document}

\section*{Задача 1.7}

\textbf{Условие:} Алфавит $\{\#,a,b\}$. Найти самое левое вхождение
подстроки $aababab$ и стереть все символы до этого вхождения.
Если подстрока не найдена, лента должна быть пуста.

\section*{Программа}

\footnotesize
\begin{verbatim}
alphabet = [#, a, b]
tapes = [ right ]
START:
    # -> #R Search_0
Search_0:
    a -> #R Search_a
    b -> #R Search_0
    # -> #H STOP
Search_a:
    a -> #R Search_aa
    b -> #R Search_0
    # -> #H STOP
Search_aa:
    a -> #R Search_aa
    b -> #R Search_aab
    # -> #H STOP
Search_aab:
    a -> #R Search_aaba
    b -> #R Search_0
    # -> #H STOP
Search_aaba:
    a -> #R Search_aa
    b -> #R Search_aabab
    # -> #H STOP
Search_aabab:
    a -> #R Search_aababa
    b -> #R Search_0
    # -> #H STOP
Search_aababa:
    a -> #R Search_aa
    b -> bL Restore_6
    # -> #H STOP
Restore_6:
    # -> aL Restore_5
    a -> aL Restore_5
    b -> aL Restore_5
Restore_5:
    # -> bL Restore_4
    a -> bL Restore_4
    b -> bL Restore_4
Restore_4:
    # -> aL Restore_3
    a -> aL Restore_3
    b -> aL Restore_3
\end{verbatim}
\newpage
\begin{verbatim}
Restore_3:
    # -> bL Restore_2
    a -> bL Restore_2
    b -> bL Restore_2
Restore_2:
    # -> aL Restore_1
    a -> aL Restore_1
    b -> aL Restore_1
Restore_1:
    # -> aH STOP
    a -> aH STOP
    b -> aH STOP
\end{verbatim}
\normalsize

\subsection*{Сложность реализации в худшем}

Пусть $n = |w|$, без $\#$.

\paragraph{Если $aababab\not\in w$.}
Переход из $START$ занимает один шаг. Затем каждый из $n$ символов
прочитывается, заменяется на $\#$, и головка перемещается вправо.
После этого выполняется переход на правой границе в $STOP$. Поэтому
\[
T_{not~in}(n)=1+n+1=n+2
\]

\paragraph{Если $aababab\in w$.}
Пусть $s$ - индекс начала первого вхождения $aababab$.
До чтения последнего символа найденной подстроки выполняется
$s+5$ шагов по символам $1,\ldots,s+5$.
Плюс начальных шаг из $START$, один переход
по последнему символу $b$ в $Restore_6$ и шесть переходов
которые восстанавливают $aababa$. Значит,
\[
T_{in}(n,s)
=1+(s+5)+1+6=s+13
\]
Так как $|aababab|=7$, то $s\leq n-6$.
Ну и тогда, при $n\geq7$
\[
T_{in}(n,s)=s+13\leq n+7
\]
И тогда самый сложный случай - $*^{n-7}aababab$, где 
в $*^{n-7}$ нет вхождения $aababab$.

Значит в худшем случае сложность реализации программы:
\[
T_{\max}(n)=
\begin{cases}
n+2, & 0\leq n<7\\
n+7, & n\geq7
\end{cases}
\]

\subsection*{Оптимальность}

Рассмотрим произвольную Машину Тьюринга $A$ с одной лентой и
входными данными $w=\#\{a,b\}^n\#$, которая корректно
решает задачу.

\paragraph{Если $0\leq n<7$.}
Вхождения $aababab$ в $w$ понятно что нет, поэтому все
входные символы должны быть стёрты. Кроме того, машина
должна обнаружить конец входа. Ну и так как машина не знает
заранее, что длина слова $<7$, то она должна прочитать все $n$
символов, чтобы убедиться в отсутствии искомой подстроки.
Тогда в худшем случае требуется один стартовый переход,
$n$ переходов по символам входа и один переход на правой границе.
Значит, для любой такой машины $A$:
\[
T_A^{\max}(n)\geq n+2, 0\leq n<7
\]

\paragraph{Если $n\geq7$.}
Рассмотрим слова длины $n$:
\[
x_n=b^{n-7}aababab,
y_n=b^{n-7}aababaa
\]
$x_n$ содержит подстроку $aababab$, а $y_n$ - нет. 
Они различаются только последним символом - на
последнем месте. До первого перехода, 
выполняемого при головке на ячейке $n$,
вычисления на этих словах очевидно из жетерминированности
совпадают. В частности, содержимое $n-6$-ого символа $c$
в этот момент одинаково.

Значение $c$ - это $a$ или $\#$ (для $b$ - не оптимально,
так как его надо заменить на $a$ или $\#$. Просто не будем
его рассматривать). Поэтому хотя бы на одном
из двух входов машина должна изменить содержимое ячейки $n-6$ уже
после того, как прочитала последний символ в ячейке $n$.

Когда головка находится в ячейке $n$, для записи в ячейку $n-6$ ей
необходимо сделать как минимум шесть перемещений влево, а затем ещё
один переход, изменяющий символ в ячейке $n-6$. Значит
\[
T_A^{\max}(n)\geq \max\{T_A(x_n),T_A(y_n)\}\geq n+7, n\geq7
\]

А значит, для любой такой машины $A$:
\[
T_A^{\max}(n)\geq
\begin{cases}
n+2, & 0\leq n<7\\
n+7, & n\geq7
\end{cases}
\]

Ну а реализация достигает этих нижних границ. Поэтому
\[
T_{\max}(n) = \min_A T_A^{\max}(n)=
\begin{cases}
n+2, & 0\leq n<7\\
n+7, & n\geq7
\end{cases}
\]

\end{document}